% RESUMO--------------------------------------------------------------------------------

\begin{resumo}[RESUMO]
\begin{SingleSpacing}

Este trabalho teve como propósito desenvolver e propor uma solução tecnológica, concretizada em um aplicativo móvel, para o registro de atividades pesqueiras, abrangendo a localização dos pontos de captura, as espécies de peixe e os tipos de iscas utilizadas. A pesquisa foi motivada pela dificuldade encontrada na rotina dos pescadores em localizar áreas de pesca que conciliem produtividade, segurança (baixo risco de acidentes) e acessibilidade. O ambiente digital FishSpot foi criado para suprir essa lacuna, oferecendo funcionalidades de apoio direto ao pescador. Como resultado, obteve-se um protótipo funcional capaz de centralizar dados cruciais para a atividade pesqueira. Conclui-se que a aplicação representa uma ferramenta com potencial significativo para promover o engajamento e a cooperação mútua na comunidade pesqueira, otimizando práticas e a gestão de recursos. 

\hfill \break
\textbf{Palavras-chave}: Pesca. Plataforma Digital. Aplicação Móvel. Software, Desenvolvimento.

\end{SingleSpacing}
\end{resumo}

